% !TeX program = pdflatex
%
% Documentation of the bill-of-materials / cost-reporting subsystem
% (cost_report.py) of the ZeroLevelProduct CAD assembly toolchain.
%
% Revised 2026-07-22: describes the simplified reporter, from which the
% optional fluid/coolant costing path has been removed. See
% claude_mds/cost_report_simplification_2026-07-22.md for the change log.
%
% extarticle (from the extsizes bundle) is a drop-in clone of article that
% accepts sizes below 10pt. Every other size in this file is RELATIVE
% (\footnotesize, \scriptsize, ...), so the base size set here scales the whole
% document -- body, tables, code listings and diagram labels alike.
% If extsizes is not installed, revert to:  \documentclass[10pt,a4paper]{article}
\documentclass[9pt,a4paper]{extarticle}

\usepackage[T1]{fontenc}
\usepackage[utf8]{inputenc}
\usepackage[margin=2.5cm]{geometry}
\usepackage{amsmath}
\usepackage{amssymb}
\usepackage{booktabs}
\usepackage{longtable}
\usepackage{array}
\usepackage{xcolor}
\usepackage{listings}
\usepackage{tikz}
\usetikzlibrary{arrows.meta,positioning,fit,backgrounds}
\usepackage[hidelinks]{hyperref}

\definecolor{codebg}{gray}{0.96}
\definecolor{codekw}{RGB}{0,80,160}
\definecolor{codecm}{RGB}{100,120,100}
\definecolor{codestr}{RGB}{160,60,60}
\definecolor{warnbg}{RGB}{255,247,230}
\definecolor{warnrule}{RGB}{200,140,40}

\lstset{
  basicstyle=\ttfamily\scriptsize,
  backgroundcolor=\color{codebg},
  keywordstyle=\color{codekw}\bfseries,
  commentstyle=\color{codecm}\itshape,
  stringstyle=\color{codestr},
  showstringspaces=false,
  breaklines=true,
  frame=single,
  rulecolor=\color{gray!40},
  framesep=4pt,
  xleftmargin=6pt,
  xrightmargin=2pt,
  literate={³}{{\textsuperscript{3}}}1 {·}{{$\cdot$}}1 {→}{{$\rightarrow$}}1
           {ρ}{{$\rho$}}1 {Σ}{{$\Sigma$}}1 {—}{{---}}1 {’}{{'}}1,
}
\lstdefinestyle{py}{language=Python}

\newcommand{\code}[1]{\texttt{\small #1}}
\newcommand{\file}[1]{\texttt{\small #1}}

% A boxed caveat, for the findings that a reader must not miss.
% Plain LaTeX (lrbox + minipage + fcolorbox) rather than a tikz node, so that it
% survives page breaks in the surrounding text and needs no extra package.
\newsavebox{\caveatbox}
\newenvironment{caveat}[1]%
  {\par\medskip\noindent
   \begin{lrbox}{\caveatbox}%
   \begin{minipage}{\dimexpr\linewidth-2\fboxsep-2\fboxrule\relax}%
   \textbf{#1}\par\medskip}%
  {\end{minipage}%
   \end{lrbox}%
   \fcolorbox{warnrule}{warnbg}{\usebox{\caveatbox}}\par\medskip}

\title{\textbf{Cost Reporting in the Zero-Level Product Toolchain}\\[0.4em]
       \large Bill of materials, mass accounting and pricing from CAD geometry}
\author{Documentation of \file{cost\_report.py}}
\date{\today}

\begin{document}
\maketitle

\begin{abstract}
\noindent
This document describes how the cost-reporting subsystem of the zero-level
reactor product model works. The subsystem takes the \emph{same} component
specification list that drives the CAD assembly, measures the true CAD volume of
every \emph{solid} material zone of every component, converts those volumes into
masses using densities derived from the neutronics material library, prices the
masses per kilogram, and writes a two-sheet spreadsheet bill of materials. The
document covers the data flow, the physical conversions and their unit handling,
the aggregation and reporting rules, the validation and error behaviour, and
what is deliberately excluded from the bill. A worked example on the ESFR-SMR
reference assembly is included, together with an audit of the known
inconsistencies in the current model --- chief among them a fuel composition
that does not match the geometry it is applied to, and which dominates the
reported total.
\end{abstract}

\tableofcontents
\bigskip

% =====================================================================
\section{Purpose and design intent}
% =====================================================================

The geometry model already knows two things that, taken together, are enough to
cost a reactor: \emph{how big every part is} (CAD volume) and \emph{what every
part is made of} (the material assignment each specification carries for the
neutronics homogeniser). The cost reporter exploits exactly this coincidence: it
adds \emph{no new geometry and no new material data}. It re-uses the existing
zone model and the existing material library, and adds only two pieces of
information of its own:

\begin{enumerate}
  \item a price list, \code{COSTS}, in currency per kilogram, and
  \item the meaning of one model length unit, so that a volume can become a mass.
\end{enumerate}

The central design decision is that materials are priced \textbf{per kilogram,
not per unit volume}, because that is how metals and nuclear fuel are actually
quoted commercially. Pricing per kilogram forces the pipeline to compute mass,
which in turn forces it to know both a density and a length unit --- and the
module is explicit that this is the one place in the toolchain that is
\emph{not} unit-agnostic.

A second decision, equally deliberate, is one of \emph{scope}: the report bills
\textbf{what the components are made of, not what they contain}. This is
developed in Section~\ref{sec:scope}.

% =====================================================================
\section{Where it sits in the pipeline}
% =====================================================================

\subsection{Two entry points}

There are two ways to obtain a report, and they are the same code path:

\begin{lstlisting}[style=py]
# 1. as a side effect of building the assembly
assemble_objects(SPECS, export_path="output/model.step", units="m",
                 costs_report=True)

# 2. standalone, without building or exporting an assembly
from cost_report import write_cost_report
write_cost_report(SPECS, units="m")
\end{lstlisting}

Both consume the identical specification list that
\code{assemble\_objects()} and \code{homogenise\_assembly()} consume, so a cost
report never requires a separate input file, a separate parts list, or any
duplicated bookkeeping. The bill of materials cannot drift away from the
geometry, because it \emph{is} the geometry.

The \code{costs\_report} argument accepts a boolean or a dictionary of options
forwarded to \code{write\_cost\_report()}, so scenario studies need not leave the
call site:

\begin{lstlisting}[style=py]
costs_report=True                                # defaults
costs_report={"costs": MY_PRICES, "currency": "USD"}
costs_report={"materials": MY_LIB}
costs_report={"out_path": "bom.xlsx"}
costs_report={}                                  # == True
costs_report=False                               # no report
\end{lstlisting}

When invoked through \code{assemble\_objects}, the \code{units} argument becomes
mandatory even if no STEP file is exported, and this is checked \emph{before}
any geometry is built, so that a specification list which cannot produce the
report it asked for fails in seconds rather than after a multi-minute build:

\begin{lstlisting}[style=py]
if report_opts is not None and units is None:
    raise ValueError(
        "assemble_objects: 'units' is required when costs_report=True - the "
        "report prices materials per kilogram, so it has to know what one "
        "model length unit is to turn volumes into masses. ...")
\end{lstlisting}

\subsection{Data flow}

\begin{figure}[h]
\centering
\begin{tikzpicture}[
  node distance=6mm,
  box/.style={draw, rounded corners=2pt, align=center, font=\scriptsize,
              minimum height=8mm, inner sep=4pt, fill=white},
  src/.style={box, fill=blue!6},
  out/.style={box, fill=green!8},
  arr/.style={-{Latex[length=2mm]}, gray!70, thick}
]
\node[src] (specs)  {\textbf{SPECS}\\component dicts\\(geometry + materials)};
\node[box, right=14mm of specs] (zones) {\code{zone\_model(spec)}\\
      \code{MATERIAL\_ZONES} $\rightarrow$ per-component\\
      zones; else \code{generic\_zones}};
\node[box, right=14mm of zones] (vol) {\code{material\_volumes()}\\
      \emph{solid} zones only\\ (id, type, zone, mat, $V$)};

\node[src, below=13mm of specs] (matlib) {\file{materials.py}\\\code{MATERIALS}};
\node[box, right=14mm of matlib] (rho) {\code{material\_density()}\\
      $\rho$ in g/cm$^3$};
\node[box, right=14mm of rho] (mass) {mass $=V\,k_u\,\rho$\\
      $k_u$ from \code{units}};

\node[src, below=13mm of matlib] (costs) {\code{COSTS}\\currency per kg};
\node[box, right=14mm of costs] (cost) {cost $=c\cdot$ mass};
\node[out, right=14mm of cost] (xlsx) {\textbf{.xlsx report}\\
      ``By material''\\``By component''};

\draw[arr] (specs) -- (zones);
\draw[arr] (zones) -- (vol);
\draw[arr] (matlib) -- (rho);
\draw[arr] (rho) -- (mass);
\draw[arr] (costs) -- (cost);
\draw[arr] (cost) -- (xlsx);
\draw[arr] (vol.south) -- ++(0,-4mm) -| (mass.north);
\draw[arr] (mass.south) -- ++(0,-4mm) -| (cost.north);
\end{tikzpicture}
\caption{Data flow of the cost reporter. Blue boxes are inputs the user
supplies or configures; the green box is the produced artefact. Nothing in the
chain introduces geometry or composition data that the CAD/neutronics model does
not already contain.}
\end{figure}

% =====================================================================
\section{Step 1 --- enumerating and measuring material zones}
% =====================================================================

\subsection{The zone model}

A \emph{zone} is a named region of a component made of a single material: a
steel shell, a bore, a cavity, a pool side. Zones are declared per component
type in \file{component\_material\_zones.py} through the \code{MATERIAL\_ZONES}
registry, and each zone function returns a list of
\code{Zone(name, role, volume)} records whose volumes are measured off real CAD
solids in the component's local frame.

Registration is optional. \code{homogenise.zone\_model(spec)} dispatches:

\begin{lstlisting}[style=py]
def zone_model(spec):
    """Registered zone declaration if there is one, generic single zone otherwise."""
    obj_type = spec.get("obj_type")
    if obj_type in MATERIAL_ZONES:
        return MATERIAL_ZONES[obj_type](spec)
    return generic_zones(spec)
\end{lstlisting}

\code{generic\_zones} builds the component through the very same
\code{build\_solid()} the assembler uses and treats the whole solid as one zone.
This is what makes the cost reporter work uniformly across all three build paths
(pre-made component, 2D profile + operation, CadQuery 3D primitive) without any
component-specific code in the reporter itself.

Registered zone functions are used where one material cannot describe a
component. The intermediate heat exchanger is the motivating case: it returns
three zones, obtained by genuine boolean geometry rather than by formula
(Table~\ref{tab:zones}).

\begin{table}[h]
\centering
\footnotesize
\begin{tabular}{@{}llll@{}}
\toprule
\textbf{Component type} & \textbf{Zones} & \textbf{Role} & \textbf{Billed?} \\
\midrule
\code{ihx}                  & \code{structure}   & solid & yes \\
                            & \code{tube\_side}  & fluid & no  \\
                            & \code{shell\_side} & fluid & no  \\
\code{diagrid}              & \code{shell}       & solid & yes \\
                            & \code{cavity}      & fluid & no  \\
\code{reactor\_core}        & \code{core}        & solid & yes \\
\code{reactor\_vessel}      & \code{structure}   & solid & yes \\
\code{reactor\_top\_plate}  & \code{structure}   & solid & yes \\
\code{primary\_pump}, \code{strongback},    & \code{structure} & solid & yes \\
\code{above\_core\_structure}, \code{redan} & \code{structure} & solid & yes \\
\midrule
\emph{anything unregistered} & single zone (\code{generic\_zones}) & solid & yes \\
\bottomrule
\end{tabular}
\caption{The \code{MATERIAL\_ZONES} registry. The \code{role} field is
informational to the homogeniser --- atom conservation treats both roles
identically --- but it is what the cost reporter filters on.}
\label{tab:zones}
\end{table}

For the IHX, the tube side is built as the union of the plenum interiors, the
tube bores, the central pipe, the riser and the lateral pipe, minus the steel
structure; the shell side is the bundle-shell bore minus the structure minus the
tube side. Those sodium inventories are therefore genuinely \emph{measured} and
remain available to the homogeniser --- the cost reporter simply does not bill
them.

\subsection{Assembling the row list}

\code{material\_volumes()} produces one row per \emph{solid} component zone:

\begin{lstlisting}[style=py]
rows: List[Tuple[str, str, str, str, float]] = []
for spec in specs:
    name   = spec.get("obj_id", spec.get("obj_type"))
    zones  = zone_model(spec)["zones"]
    assign = _zone_materials(spec, {z.name for z in zones})
    for z in zones:
        if z.name not in assign:
            raise KeyError(f"{name}: zone {z.name!r} has no material assigned. "
                           f"Assigned: {sorted(assign)}.")
        if z.role == "fluid":
            continue              # what it holds, not what it is made of
        rows.append((spec.get("obj_id", "?"), spec.get("obj_type", "?"),
                     z.name, assign[z.name], z.volume))
\end{lstlisting}

Three details are worth noting.

\paragraph{The material binding is checked for every zone, billed or not.}
The \code{KeyError} is raised \emph{before} the \code{fluid} test, so a
specification that is incomplete for the homogeniser is also rejected here. Had
the order been reversed, a missing material on a fluid zone would have passed the
cost report and failed only later, in neutronics.

\paragraph{Nothing is ever measured off an unresolved specification.}
\code{resolve()} and \code{auto\_generate\_topplate\_holes()} run
unconditionally, even when the caller has already run them --- as
\code{assemble\_objects} has:

\begin{lstlisting}[style=py]
is_premade = [d.get("obj_type") in PREMADE_BUILDERS for d in specs]
if any(is_premade):
    resolved = iter(resolve([d for d, p in zip(specs, is_premade) if p]))
    specs = [next(resolved) if p else d for d, p in zip(specs, is_premade)]
auto_generate_topplate_holes(specs)
\end{lstlisting}

Both passes are idempotent: they assign through \code{setdefault}, and the hole
pass re-reads its own previous output as user holes and adds nothing. Running
them twice costs a few milliseconds of dictionary work and no geometry;
\emph{not} running them once would silently price a component at its
pre-resolver dimensions. The top-plate pass matters physically --- without it
the plate is measured as an unperforated slab and its steel is over-counted by
the volume of every penetration (7.09~m$^3$, or 22\,\%, in the reference case).

Only pre-made components are sent through \code{resolve()}, because
\code{resolve()} deep-copies its input, which would downcast the live OCCT
shapes carried by profile-based specifications. The copy has the useful side
effect of keeping this pass from writing resolved keys back into the caller's
own dictionaries.

\paragraph{Cost.} Zone functions measure their own solids, which means each
component is rebuilt for the report. Requesting a report therefore costs roughly
one extra assembly build.

\subsection{Binding zones to materials}

\code{\_zone\_materials()} resolves the zone $\rightarrow$ material map from
either accepted form: a \code{"materials"} block for multi-zone components, or
the single \code{"material"} key for one-zone components.

\begin{lstlisting}[style=py]
if "materials" in spec:
    assign = dict(spec["materials"])
    assign.pop(FILLER_KEY, None)          # "_filler" is not a zone
    return assign
if "material" in spec:
    if len(declared) != 1:
        raise KeyError("... needs a 'materials' block; the single 'material' "
                       "key is only valid for one-zone components.")
    return {next(iter(declared)): spec["material"]}
raise KeyError("... no material assignment ...")
\end{lstlisting}

This is deliberately \emph{looser} than the homogeniser's equivalent
(\code{homogenise.\_assignment}), which additionally requires the single zone to
be named by the reserved solid-zone constant. Here any one-zone component may
use the short form whatever its zone happens to be called
(\code{"structure"}, \code{"core"}, \dots), because pricing has no stake in the
zone-naming contract that atom conservation depends on. The reactor vessel in
the reference example exploits exactly this: it declares a zone named
\code{structure} but carries only \code{"material": "ss316"}.

Conversely, a declared zone with no material raises \code{KeyError}.

% =====================================================================
\section{Step 2 --- density from the material library}
% =====================================================================

Densities are \emph{not} extra data the user must supply for costing. They are
derived from the same \code{MATERIALS} library the neutronics workflow uses,
which stores materials in one of two forms.

\subsection{Materials given as a mass density}

Entries of kind \code{"mass\_density"} state $\rho$ outright in g/cm$^3$; it is
read directly. The sodium coolants are given this way, because their density is
a strong function of temperature and is better stated explicitly for the
primary- and secondary-side conditions than derived:

\begin{lstlisting}[style=py]
"na_primary": {"kind": "mass_density", "density": 0.83,   # g/cm3
               "nuclides": {"Na23": {"ao": 1.0, "M": 22.99}}},
\end{lstlisting}

\subsection{Materials given as number densities}
\label{sec:numdens}

Entries of kind \code{"number\_density"} give, per nuclide, an atom density
$n_i$ in atoms/(b$\cdot$cm) --- the unit neutronics codes want. The mass density
follows from

\begin{equation}
  \rho \;=\; \frac{10^{24}}{N_A}\sum_i n_i A_i
  \qquad
  \left[\frac{\text{g}}{\text{cm}^3}\right],
  \label{eq:rho}
\end{equation}

where $10^{24}$ converts atoms/(b$\cdot$cm) to atoms/cm$^3$, $N_A =
6.022\,140\,76\times10^{23}\,\text{mol}^{-1}$ is imported from
\file{homogenise.py} so that both workflows share one constant, and $A_i$ is the
mass number of nuclide $i$, \emph{parsed out of the nuclide's own name} by
regular expression (\code{Fe56} $\rightarrow$ 56):

\begin{lstlisting}[style=py]
total = 0.0
for nuc, n in mat["nuclides"].items():
    a = re.search(r"(\d+)", nuc)
    if a is None:
        raise ValueError(f"cannot read a mass number off nuclide name {nuc!r} ...")
    total += float(n) * float(a.group(1))
return total * 1e24 / N_A
\end{lstlisting}

Using the mass number in place of the molar mass is an approximation at the
level of the nuclear mass defect, $\sim 0.1\,\%$. That is orders of magnitude
finer than the uncertainty on any price list, and it spares the material library
from carrying molar masses it does not otherwise need. Where a name carries no
digits the code refuses to guess and directs the user to supply the material as
a \code{mass\_density} entry instead.

\paragraph{Worked check.} For \code{ss316}, with
$n_{\mathrm{Fe56}}=5.90\times10^{-2}$, $n_{\mathrm{Cr52}}=1.60\times10^{-2}$,
$n_{\mathrm{Ni58}}=8.00\times10^{-3}$, $n_{\mathrm{Mo98}}=1.20\times10^{-3}$:
\begin{align*}
  \sum_i n_i A_i &= 3.304 + 0.832 + 0.464 + 0.1176 = 4.7176 \\
  \rho &= 4.7176 \times \frac{10^{24}}{6.02214076\times10^{23}}
        = 7.834\ \text{g/cm}^3 .
\end{align*}

\begin{caveat}{Caveat --- natural-element nuclide names}
The regular expression matches the \emph{first} run of digits, so the OpenMC
convention for a natural element, \code{C0} or \code{Fe0}, parses as $A = 0$ and
contributes \emph{zero} mass with no diagnostic. The function refuses a name with
no digits at all, but silently accepts the one pattern that means ``natural
element''. No material in the current library is affected; a library that adds
one would be silently too light. See Section~\ref{sec:audit}.
\end{caveat}

% =====================================================================
\section{Step 3 --- units: from model volume to kilograms}
% =====================================================================

Zone volumes come out of CAD in the model's own length unit cubed; densities are
in g/cm$^3$. Turning one into the other requires knowing what one model unit
physically \emph{is}, which is precisely why \code{units} is a required
argument and not an inferred one. This is the module's only unit knowledge:

\begin{lstlisting}[style=py]
_CM_PER_UNIT = {"m": 100.0, "mm": 0.1, "cm": 1.0, "um": 1e-4, "km": 1e5,
                "inch": 2.54, "ft": 30.48}
...
kg_per_unit3 = _CM_PER_UNIT[unit_key] ** 3 / 1000.0
\end{lstlisting}

Writing $\ell_u$ for the centimetres in one model unit, the conversion factor is

\begin{equation}
  k_u \;=\; \frac{\ell_u^{\,3}}{1000},
  \qquad\text{so}\qquad
  m\,[\text{kg}] \;=\; V\,[\text{unit}^3]\;\cdot\;k_u\;\cdot\;\rho\,[\text{g/cm}^3],
  \label{eq:mass}
\end{equation}

where $\ell_u^{\,3}$ converts unit$^3$ to cm$^3$ and the factor $1000$ converts
grams to kilograms. The keys of \code{\_CM\_PER\_UNIT} deliberately match
\code{assemble.\_STEP\_UNITS}, so a specification's \code{units=} means the same
thing in the STEP header and in the bill of materials. \textbf{Nothing is
rescaled} by the reporter: the unit is used only inside
Equation~\eqref{eq:mass}, and the geometry itself remains unit-agnostic. An
unrecognised unit is rejected with the accepted list.

The self-test at the bottom of the module asserts the invariant that makes this
safe --- one cubic metre of steel must weigh exactly what the $10^9$ cubic
millimetres describing the very same block weigh:

\begin{lstlisting}[style=py]
kg = lambda v, u: v * _CM_PER_UNIT[u] ** 3 / 1000.0 * rho
assert abs(kg(1.0, "m") - kg(1e9, "mm")) < 1e-6 * kg(1.0, "m")
\end{lstlisting}

The same self-test checks the derived densities against reference values
(\code{ss316} $7.83$, \code{ss304} $7.72$, \code{core\_smear} $10.48$,
\code{na\_primary} $0.83$ g/cm$^3$) and runs in pure Python, with no CadQuery
dependency.

% =====================================================================
\section{Step 4 --- pricing and aggregation}
\label{sec:pricing}
% =====================================================================

\subsection{The price list}

\code{COSTS} maps a material name to a price in currency per kilogram, and
\code{CURRENCY} is a pure \emph{label} --- no conversion of any kind is
performed, so the label must match whatever the numbers are actually in. Its
only purpose is to ensure that no column in the report is a bare number whose
unit the reader has to guess.

\begin{table}[h]
\centering
\footnotesize
\begin{tabular}{@{}lr@{}}
\toprule
\textbf{Material} & \textbf{Price (EUR/kg)} \\
\midrule
\code{ss316}        & 8 \\
\code{ss304}        & 5 \\
\code{core\_smear}  & 10\,000 \\
\bottomrule
\end{tabular}
\caption{The default \code{COSTS} table (placeholder figures, intended to be
replaced with project-specific quotations). The coolants carry no entry: fluid
zones are never billed, so a price for one would be dead data.}
\end{table}

\paragraph{A caveat on fuel.} Nuclear fuel is normally quoted per kilogram of
\emph{heavy metal}, whereas the mass computed here is that of the whole
compound. For \code{core\_smear} the heavy-metal mass fraction is
\begin{equation}
  f_{\mathrm{HM}}
  = \frac{\sum_{i \in \mathrm{HM}} n_i A_i}{\sum_i n_i A_i}
  = \frac{5.569}{6.312} = 0.8823 ,
\end{equation}
the balance being oxygen. Since the reported cost is $c\,M_{\text{compound}}$
whereas a quotation gives $P_{\mathrm{HM}} M_{\mathrm{HM}}$, and
$M_{\mathrm{HM}} = f_{\mathrm{HM}} M_{\text{compound}}$, the value to enter in
\code{COSTS} is
\begin{equation}
  c \;=\; f_{\mathrm{HM}}\, P_{\mathrm{HM}} \;=\; 0.8823\,P_{\mathrm{HM}} ,
  \label{eq:fuelprice}
\end{equation}
that is, a kgHM price scaled \textbf{down} by the heavy-metal fraction.

\begin{caveat}{Caveat --- the docstring states this backwards}
The docstring above \code{COSTS} in \file{cost\_report.py} instructs the reader
to ``scale a kgHM price \emph{up} by that fraction''. Equation~\eqref{eq:fuelprice}
is the correct direction. Because fuel is $99.9\,\%$ of the reference total
(Section~\ref{sec:worked}), following the docstring literally puts a factor
$1/f_{\mathrm{HM}}^2 \approx 1.28$ error on essentially the whole bill.
\end{caveat}

\subsection{The pricing functions}

Costing is entirely linear in volume:

\begin{lstlisting}[style=py]
def mass_of(mat, vol):        # kg
    return vol * kg_per_unit3 * density[mat]

def cost_of(mat, vol):        # currency, or None if unpriced
    return None if mat not in costs else costs[mat] * mass_of(mat, vol)

def cost_cell(cost):          # None must not be written as a blank
    return UNPRICED if cost is None else cost

def millions(cost):
    return UNPRICED if cost is None else cost / 1e6
\end{lstlisting}

so that for a solid zone $z$ of material $m(z)$ and volume $V_z$,

\begin{equation}
  C \;=\; \sum_{z\ \mathrm{solid}} c_{m(z)}\; k_u\; \rho_{m(z)}\; V_z ,
\end{equation}

with $c_m$ the price per kilogram. Because the relation is linear, summing costs
over components and summing costs over per-material total volumes give
identical answers --- the two sheets are guaranteed consistent by construction,
not by a reconciliation step.

\subsection{Unpriced materials}

A material present in the geometry but absent from \code{COSTS} is \emph{not}
silently dropped. Its cost cells read the sentinel string \code{UNPRICED}, its
mass still counts towards the mass total, and the \code{TOTAL} label is changed
to \code{TOTAL (excl.\ UNPRICED)} with the offending materials named in a note
beneath the table:

\begin{lstlisting}[style=py]
grand_total = sum(c for c in (cost_of(m, v) for m, v in totals.items())
                  if c is not None)
ws.append([f"TOTAL (excl. {UNPRICED})" if unpriced else "TOTAL", ...])
\end{lstlisting}

A blank cell was rejected deliberately: in a spreadsheet a blank is
indistinguishable from zero to anyone summing a column, and it would quietly
shrink the bill. A word cannot be summed by accident. The mass total, by
contrast, covers everything --- mass is never unknown, only price is.

\subsection{Why cost appears twice}

Every cost is reported both in currency units and in millions of currency units.
Reactor bills of materials run to eight and nine figures, at which point the raw
number is unreadable, while the millions figure is the one that ends up in the
write-up. Reporting both avoids a manual division --- and the transcription
error that goes with it.

% =====================================================================
\section{The output workbook}
% =====================================================================

The report is an \code{.xlsx} workbook written with \code{openpyxl}, containing
two sheets. Every column header carries its unit: volume headers are stamped
with whichever \code{units} was passed, money headers with \code{CURRENCY}.
Header rows are bold and column widths are fitted to their content, capped at 40
characters so that the \code{UNPRICED} footnote --- one long cell in column A ---
does not stretch the material column across the screen.

\subsection{Sheet 1 --- ``By material''}

One row per distinct material, plus a \code{TOTAL} row:

\begin{center}\footnotesize
\code{Material} $\mid$ \code{Total volume (unit³)} $\mid$ \code{Density (g/cm³)}
$\mid$ \code{Total mass (kg)} $\mid$ \code{Cost per kg (CUR/kg)} $\mid$
\code{Total cost (CUR)} $\mid$ \code{Total cost (M CUR)}
\end{center}

Materials are listed alphabetically so that successive reports diff cleanly. In
the \code{TOTAL} row the density and unit-price cells are left empty on purpose:
an average of densities weighted by nothing, or of prices weighted by nothing,
would be a meaningless number.

The sheet closes with a statement of its own scope:

\begin{center}\footnotesize\ttfamily
Scope: what the components are made of. Fluid zones and coolant are NOT included.
\end{center}

This is not decoration. A reader opening the file cold cannot otherwise tell
whether the absence of a coolant line means the model has no coolant or means
the report left it out.

\subsection{Sheet 2 --- ``By component''}

One row per billed component zone, in specification order, which keeps the sheet
readable as a walk through the assembly:

\begin{center}\footnotesize
\code{Component} $\mid$ \code{Type} $\mid$ \code{Zone} $\mid$ \code{Material}
$\mid$ \code{Volume (unit³)} $\mid$ \code{Mass (kg)} $\mid$ \code{Cost (CUR)}
$\mid$ \code{Cost (M CUR)}
\end{center}

This sheet is the audit trail for the first: every figure in ``By material'' is
the sum of a subset of rows here, and each row is attributable to a named
component and a named zone of that component.

\subsection{Where the file is written}

If no \code{out\_path} is given, the report auto-names itself after
\emph{the script that asked for it}:

\begin{center}
\file{<dir of that script>/costs\_reports/<stem>\_costs\_reports\_<YYYYmmdd\_HHMMSS>.xlsx}
\end{center}

\code{caller\_report\_path()} walks the call stack outwards, past
\file{cost\_report.py} itself and past anything named in its \code{skip}
argument, and takes the first frame that belongs to user code:

\begin{lstlisting}[style=py]
ignore = {os.path.abspath(f) for f in (__file__, *skip)}
frame = next((f for f in inspect.stack()[1:]
              if os.path.abspath(f.filename) not in ignore), None)
\end{lstlisting}

\code{assemble\_objects} passes \code{skip=(\_\_file\_\_,)} so that the name
chosen is the user's example script and not \file{assemble.py}. Anchoring the
output to \emph{the script's own directory} rather than to the working directory
means the report lands in the same place whether the example is run from the
project root or from \file{examples\_reactor/}. In a REPL, notebook or
\code{exec()} context, where no source file exists, the reporter falls back to
the stem \file{assembly} in the current working directory. The timestamp makes
successive runs non-destructive, and the parent directory is created if needed.

% =====================================================================
\section{Validation, errors and warnings}
% =====================================================================

The reporter distinguishes sharply between conditions that make the report
\emph{wrong} (fatal) and conditions that make it \emph{incomplete} (a warning).

\begin{table}[h]
\centering
\footnotesize
\renewcommand{\arraystretch}{1.2}
\begin{tabular}{@{}p{5.0cm}p{2.4cm}p{7.0cm}@{}}
\toprule
\textbf{Condition} & \textbf{Behaviour} & \textbf{Rationale} \\
\midrule
\code{costs\_report=True} with no \code{units} & \code{ValueError}, before the build &
Mass, hence cost, is undefined; failing early saves a full geometry build. \\
Unrecognised \code{units} string & \code{ValueError} with the accepted list &
Silently guessing a unit would silently scale the whole bill. \\
\code{costs\_report} dict overriding \code{specs}/\code{units} & \code{ValueError} &
Those are set by \code{assemble\_objects} itself. \\
Component with no \code{material} / \code{materials} & \code{KeyError} &
An unpriced part must never be silently omitted from a bill of materials. \\
Declared zone with no material & \code{KeyError} listing what \emph{was} assigned &
Same reason, at zone granularity; checked for fluid zones too. \\
Single \code{material} key on a multi-zone component & \code{KeyError} &
The assignment is ambiguous. \\
Material used by geometry but absent from the material library & \code{KeyError} &
Its density, and therefore its mass, is unknown. \\
Nuclide name with no mass number & \code{ValueError} &
$A_i$ cannot be read; the user is told to give a \code{mass\_density} entry. \\
Material in the geometry but absent from \code{COSTS} & \code{warnings.warn}; cells read \code{UNPRICED} &
Volume, density and mass are still valid and useful; only the price is missing. \\
\bottomrule
\end{tabular}
\caption{Validation behaviour.}
\end{table}

The unpriced-material case is the only soft failure. A total obtained under that
warning is a lower bound, its label says so, and the warning names the materials
responsible.

% =====================================================================
\section{Scope: what is deliberately not counted}
\label{sec:scope}
% =====================================================================

\paragraph{Fluids, coolant and the sodium inventory.} Only zones the zone
functions mark \code{role="solid"} are billed. A bill of materials is a list of
parts to buy; the sodium a component contains is a consumable inventory, not a
part, and in a pool reactor it is most of the volume inside the vessel --- it
would swamp a column meant to tell you about hardware. Only zones
\emph{explicitly} marked \code{"fluid"} are dropped: a zone with any other role
is billed, because silently omitting something from a bill of materials is worse
than listing something the reader did not want.

An earlier revision of this module carried an opt-in \code{include\_fluids} path
that also derived, per container, the coolant standing around every component by
a boolean cut-chain. It was removed on 2026-07-22 together with the
\code{encloses} specification keys and the \code{enclosure} clause of the zone
model. Costing the coolant is a separate problem and is now firmly out of scope.

\paragraph{The homogenisation filler.} The neutronics homogeniser wraps each
component in an equivalent cylinder and fills the remaining volume with a
\code{"\_filler"} material (pool sodium rather than vacuum). That volume is a
\emph{closure remainder of a modelling construct}, not a part anyone buys. The
key is explicitly discarded by \code{\_zone\_materials()}.

\paragraph{Equivalent cylinders generally.} No homogenised cylinder takes part
in the bill. They are a deliberately generous envelope drawn round each
component, and the gap between a component and its envelope is not a quantity
anyone can buy.

\paragraph{Everything that is not material.} The report is a material bill:
mass times a price per kilogram. Fabrication, welding, machining, transport,
installation, quality assurance, contingency and every other non-material cost
element are outside its scope. It answers ``what does the material in this
design cost?'', which is the appropriate question at zero-level design.

\paragraph{Currency conversion.} None is performed. \code{CURRENCY} labels the
columns; the numbers are whatever \code{COSTS} is denominated in.

% =====================================================================
\section{Worked example: the ESFR-SMR reference assembly}
\label{sec:worked}
% =====================================================================

The reference case is \file{examples\_reactor/example\_final\_esfr\_smr\_with\_materials.py},
whose thirteen specifications (vessel, top plate, three IHXs, three primary
pumps, diagrid, core, strongback, redan and above-core structure) are assembled
in metres with

\begin{lstlisting}[style=py]
assemble_objects(SPECS, export_path=f"output/esfr_smr_full_reactor_{_TS}.step",
                 units="m", costs_report=True)
\end{lstlisting}

With \code{units="m"} the conversion factor of Equation~\eqref{eq:mass} is
$k_u = 100^3/1000 = 1000$, i.e.\ mass in kilograms is simply volume in m$^3$
times density in g/cm$^3$ times $1000$.

\subsection{Sheet 1 --- ``By material''}

\begin{table}[h]
\centering
\footnotesize
\begin{tabular}{@{}lrrrrrr@{}}
\toprule
\textbf{Material} & \textbf{Volume} & \textbf{Density} & \textbf{Mass} &
\textbf{Price} & \textbf{Cost} & \textbf{Cost} \\
 & (m$^3$) & (g/cm$^3$) & (kg) & (EUR/kg) & (EUR) & (M EUR) \\
\midrule
\code{core\_smear}   &  39.799 & 10.481 &   417\,123 & 10\,000 & 4\,171\,228\,716 & 4\,171.23 \\
\code{ss304}         &  32.177 &  7.722 &   248\,453 &       5 &      1\,242\,265 &     1.24 \\
\code{ss316}         &  71.300 &  7.834 &   558\,544 &       8 &      4\,468\,348 &     4.47 \\
\midrule
\textbf{TOTAL}       & 143.275 &  ---   & 1\,224\,119 &   ---   & 4\,176\,939\,329 & 4\,176.94 \\
\bottomrule
\end{tabular}
\caption{``By material'' sheet for the ESFR-SMR reference assembly, with the
placeholder price list. Densities are derived from \code{MATERIALS} through
Equation~\eqref{eq:rho}.}
\label{tab:bymaterial}
\end{table}

The structure of the result is worth reading rather than merely tabulating. The
hardware amounts to roughly $1\,224$ tonnes, of which $807$ tonnes are steel and
$417$ tonnes are core. Yet at these placeholder prices the core alone is
$99.87\,\%$ of the bill: fuel at $10^4$ EUR/kg against steel at $8$ EUR/kg is a
three-orders-of-magnitude ratio that no plausible steel mass can offset.

That is the kind of conclusion the report exists to make immediate. It also
means the total is only as good as the fuel model behind it, and
Section~\ref{sec:audit} shows that the fuel model is currently the weakest part
of the chain. Note the \code{TOTAL} volume cell adds steel and fuel volumes
together; the mass and cost totals are genuine sums, the volume total is not a
physical quantity.

\subsection{Sheet 2 --- ``By component''}

\begin{footnotesize}
\begin{longtable}{@{}llllrrr@{}}
\caption{``By component'' sheet for the same run. The three IHXs and three pumps
are identical instances and reproduce identical figures, which is a useful
sanity check on the zone measurement. The diagrid contributes only its steel
shell; its sodium-filled cavity is measured but not billed.}
\label{tab:bycomponent}\\
\toprule
\textbf{Component} & \textbf{Type} & \textbf{Zone} & \textbf{Material} &
\textbf{Vol.\ (m$^3$)} & \textbf{Mass (kg)} & \textbf{Cost (EUR)} \\
\midrule
\endfirsthead
\toprule
\textbf{Component} & \textbf{Type} & \textbf{Zone} & \textbf{Material} &
\textbf{Vol.\ (m$^3$)} & \textbf{Mass (kg)} & \textbf{Cost (EUR)} \\
\midrule
\endhead
\code{rv}         & \code{reactor\_vessel}    & structure & ss316       & 17.763 & 139\,147 & 1\,113\,179 \\
\code{top\_plate} & \code{reactor\_top\_plate}& structure & ss304       & 32.177 & 248\,453 & 1\,242\,265 \\
\code{ihx\_1}     & \code{ihx}                & structure & ss316       &  1.522 &  11\,924 &    95\,392 \\
\code{ihx\_2}     & \code{ihx}                & structure & ss316       &  1.522 &  11\,924 &    95\,392 \\
\code{ihx\_3}     & \code{ihx}                & structure & ss316       &  1.522 &  11\,924 &    95\,392 \\
\code{pump\_1}    & \code{primary\_pump}      & structure & ss316       &  2.223 &  17\,415 &   139\,323 \\
\code{pump\_2}    & \code{primary\_pump}      & structure & ss316       &  2.223 &  17\,415 &   139\,323 \\
\code{pump\_3}    & \code{primary\_pump}      & structure & ss316       &  2.223 &  17\,415 &   139\,323 \\
\code{diagrid}    & \code{diagrid}            & shell     & ss316       &  1.491 &  11\,678 &    93\,422 \\
\code{core}       & \code{reactor\_core}      & core      & core\_smear & 39.799 & 417\,123 & 4\,171\,228\,716 \\
\code{strongback} & \code{strongback}         & structure & ss316       & 25.805 & 202\,152 & 1\,617\,217 \\
\code{redan}      & \code{redan}              & structure & ss316       &  4.311 &  33\,768 &   270\,143 \\
\code{above\_core\_structure} & \code{above\_core\_structure} & structure & ss316 & 10.695 & 83\,780 & 670\,241 \\
\bottomrule
\end{longtable}
\end{footnotesize}

\subsection{An independent check on the top plate}

The top plate is the single largest steel item, and it exercises the one pass in
the pipeline that would otherwise over-count. An unperforated slab of diameter
$10.0$~m and thickness $0.5$~m has volume $\pi \times 5.0^2 \times 0.5 =
39.270$~m$^3$; the report bills $32.177$~m$^3$, a difference of
$7.093$~m$^3$. Summing the hole volumes independently --- three IHX bores, three
pump barrels and the above-core-structure neck, each sized exactly to its
component diameter --- gives
\begin{equation*}
  3(\pi\,0.800^2) \times 0.5 \;+\; 3(\pi\,0.675^2) \times 0.5
  \;+\; (\pi\,1.1085^2) \times 0.5 \;\approx\; 7.10\ \text{m}^3 ,
\end{equation*}
which reproduces the difference. The auto-generated top-plate holes are
therefore being both cut and measured correctly.

% =====================================================================
\section{Audit: known inconsistencies}
\label{sec:audit}
% =====================================================================

The following were established by direct measurement against the reference
assembly on 2026-07-22. They are ordered by their effect on the reported total.
They fall into three kinds, and it is worth keeping them apart:

\begin{itemize}
  \item \textbf{Data that does not reflect reality} --- the core composition
        (\S\ref{sec:aud:core}), the price list (\S\ref{sec:aud:prices}) and the
        steel compositions (\S\ref{sec:aud:steel}). These dominate, and none of
        them is a software fault: the reporter prices faithfully what it is
        given.
  \item \textbf{Defects in \file{cost\_report.py}} --- the inverted fuel-price
        instruction (\S\ref{sec:aud:kghm}), silently accepted unknown zone names
        (\S\ref{sec:aud:zones}) and silently accepted $A=0$
        (\S\ref{sec:aud:natural}).
  \item \textbf{Modelling and scope limits} --- split output locations
        (\S\ref{sec:aud:paths}), no overlap guard (\S\ref{sec:aud:overlap}),
        two heavy structures worth a sanity check (\S\ref{sec:aud:heavy}) and
        the extent of the parts list itself (\S\ref{sec:aud:scope}).
\end{itemize}

Only the first item, the core, is wrong by a large factor; but it is emphatically
\emph{not} the only respect in which the model departs from reality.

\subsection{The core composition does not match the core geometry}
\label{sec:aud:core}

\code{core\_smear} is $99.87\,\%$ of the reported cost, and it rests on two
compounding inconsistencies.

\paragraph{The composition is not smeared.} Its name and its comment in
\file{materials.py} describe a homogenised SFR core, but its numbers describe a
fuel pellet:

\begin{table}[h]
\centering
\footnotesize
\begin{tabular}{@{}llp{7.2cm}@{}}
\toprule
\textbf{Check} & \textbf{Value} & \textbf{Reading} \\
\midrule
O/HM atom ratio      & $1.986$              & stoichiometric (U,Pu)O$_2$: no coolant, clad or void smeared in \\
Fe / Cr / Ni content & none                 & no cladding or structure \\
Na content           & none                 & no coolant \\
Derived density      & $10.481$ g/cm$^3$    & MOX \emph{pellet} density; a smeared SFR core is $\sim$4--5 g/cm$^3$ \\
\bottomrule
\end{tabular}
\end{table}

\paragraph{The height does not match either.} \file{materials.py} states the
blend was homogenised over $95$~cm of active fuel; the \code{CORE} specification
models a $3.910$~m cylinder, a factor $4.12$ in volume. The specification's own
comment already flags this.

\paragraph{Combined effect.} $39.799$~m$^3 \times 10.481$~g/cm$^3 = 417$~t of
oxide, hence $417 \times 0.8823 = 368$~t of heavy metal. A 1500~MWth SFR core
carries roughly $50$--$70$~t~HM, so the fissile inventory --- and with it
$99.9\,\%$ of the bill --- is over by something like a factor five to seven.

The reporter is faithfully pricing what the specification and the library tell
it. The correction belongs in \file{materials.py} and in the \code{CORE}
specification, not here. Until it is made, the steel figures in
Table~\ref{tab:bymaterial} should be read as the meaningful part of the report
and the grand total should not be quoted.

\subsection{The fuel-price docstring is inverted}
\label{sec:aud:kghm}

See the caveat in Section~\ref{sec:pricing}: the module instructs the reader to
scale a kgHM price \emph{up} by the heavy-metal fraction, where
Equation~\eqref{eq:fuelprice} requires scaling \emph{down}.

\subsection{Every price in the table is a placeholder}
\label{sec:aud:prices}

\code{COSTS} is self-declared placeholder data (``replace with your own
figures''), and no entry has been traced to a quotation. The fuel figure of
$10^4$~EUR/kg is of the right order for MOX \emph{per kilogram of heavy metal},
but it is applied per kilogram of compound --- which is exactly the error
Section~\ref{sec:aud:kghm} describes. The steel figures, $8$ and
$5$~EUR/kg, are plausible for commercial plate and almost certainly low for
nuclear-qualified forgings, where qualification, inspection and traceability
dominate the material price.

Consequently the \emph{structure} of the reported bill --- which materials
matter, and by how much --- is meaningful, while its absolute total is
indicative only.

\subsection{The steel compositions are approximate}
\label{sec:aud:steel}

\file{materials.py} flags its steels as placeholder densities, and the flag is
warranted. Each alloy is represented by a single dominant isotope per element
and omits the minor alloying additions entirely:

\begin{table}[h]
\centering
\footnotesize
\begin{tabular}{@{}lrrl@{}}
\toprule
\textbf{Material} & \textbf{Derived $\rho$} & \textbf{Literature $\rho$} & \textbf{Mass fractions as modelled} \\
\midrule
\code{ss316} & $7.834$ & $\approx 7.99$ \ ($-2.0\,\%$) & Fe $70.0$, Cr $17.6$, Ni $9.8$, Mo $2.5\,\%$ \\
\code{ss304} & $7.722$ & $\approx 8.00$ \ ($-3.5\,\%$) & Fe $72.3$, Cr $19.0$, Ni $8.7\,\%$ \\
\bottomrule
\end{tabular}
\caption{Modelled versus literature densities for the two structural steels,
in g/cm$^3$.}
\end{table}

Three approximations compound here:

\begin{enumerate}
  \item \textbf{Mono-isotopic elements.} Iron is carried as \code{Fe56} alone
        ($A=56$ against a natural $\bar{A}\approx55.85$), nickel as \code{Ni58}
        ($58$ against $58.69$), molybdenum as \code{Mo98} ($98$ against
        $95.95$). Individually these are sub-percent to $2\,\%$ effects on $A_i$
        and they partially cancel.
  \item \textbf{No minor constituents.} Manganese, silicon and carbon are absent
        --- around $2$--$3\,\%$ of a real 316 by mass.
  \item \textbf{Element balance drifts from specification.} As modelled, 316 is
        $70\,\%$ Fe and $9.8\,\%$ Ni, against roughly $65\,\%$ and $12\,\%$ for
        the real alloy.
\end{enumerate}

The net effect is that both steels come out $2$--$3.5\,\%$ light, so all
structural masses and costs in Table~\ref{tab:bymaterial} are low by that
margin. This is negligible beside the fuel term today --- but once
Section~\ref{sec:aud:core} is corrected, steel becomes essentially the whole
bill and this becomes the leading material-data error.

\subsection{Unknown zone names are silently ignored}
\label{sec:aud:zones}

A \code{"materials"} block naming a zone that does not exist is rejected by the
homogeniser and accepted by the cost reporter:

\begin{center}\footnotesize
\begin{tabular}{@{}ll@{}}
\toprule
\textbf{Consumer} & \textbf{Behaviour on a mistyped zone name} \\
\midrule
\code{homogenise()}        & raises \code{KeyError} \\
\code{material\_volumes()} & accepts silently \\
\bottomrule
\end{tabular}
\end{center}

The reporter checks that every \emph{zone} has a material, but never that every
\emph{named material} corresponds to a zone. Two consumers of one specification
therefore disagree about whether it is valid, and the permissive one is the one
that produces the money figure: a mistyped zone name means a material the user
believes they assigned is quietly not billed. Tightening
\code{material\_volumes()} to match the homogeniser would remove the asymmetry.

\subsection{Natural-element nuclide names}
\label{sec:aud:natural}

As described in the caveat in Section~\ref{sec:numdens}: \code{C0} parses as
$A=0$ and contributes zero mass silently. Latent --- no current material is
affected, though Section~\ref{sec:aud:steel} notes that adding the missing
carbon to the steels is one obvious way to trip it.

\subsection{Split output locations}
\label{sec:aud:paths}

The STEP file is written relative to the \emph{working directory}, while
\code{caller\_report\_path()} anchors the workbook to the \emph{calling script's}
directory. Running the reference example from the project root therefore places
the two artefacts of one run in \file{output/} and in
\file{examples\_reactor/costs\_reports/} respectively. Each rule is defensible
alone; together they are inconsistent.

\subsection{Overlapping components}
\label{sec:aud:overlap}

Zone volumes are summed independently, so two components occupying the same
space would each be billed in full. The reporter has no guard of its own: the
only check is the overlap warning inside \code{assemble\_objects}, which does not
run on the standalone \code{write\_cost\_report()} path and which skips pairs
whitelisted through \code{interfaces\_with}. This is currently benign --- a full
run of the reference assembly produces no overlap warnings, and no specification
in it uses \code{interfaces\_with} --- but it is an assumption, not a guarantee.

\subsection{What was verified clean}

\begin{center}\footnotesize
\begin{tabular}{@{}lp{6.6cm}@{}}
\toprule
\textbf{Check} & \textbf{Result} \\
\midrule
\code{material\_volumes()} idempotency & bit-identical rows on repeat calls and on already-resolved lists \\
Top-plate hole accounting              & reproduces the $7.093$~m$^3$ difference independently \\
Billed materials present in \code{MATERIALS} & all \\
Billed materials present in \code{COSTS}     & all; no \code{UNPRICED} rows \\
Prices with no billed material               & none \\
Unintentional component overlaps             & none \\
Density self-test                            & passes \\
Unit cancellation, 1 m$^3$ vs $10^9$ mm$^3$  & exact \\
Atom conservation (homogeniser, post-change) & closure $0$, worst relative error $1.6\times10^{-16}$ \\
\bottomrule
\end{tabular}
\end{center}

\subsection{Two figures worth a design sanity check}
\label{sec:aud:heavy}

Neither is a software fault, but both follow from the geometry model and both
are large (Table~\ref{tab:bycomponent}). The \textbf{top plate} ($32.18$~m$^3$,
$248$~t) is more steel than the
\textbf{reactor vessel} ($17.76$~m$^3$, $139$~t). The \textbf{strongback} is
billed as an essentially solid forging at $25.81$~m$^3$ and $202$~t.

Between them these two components are $450$~t of the $807$~t of steel in the
model --- $56\,\%$. If either is meant to be a lighter or more perforated
structure than the current specification produces, it will move the steel bill
more than any price assumption will.

\subsection{Scope of the parts list}
\label{sec:aud:scope}

The bill covers \emph{the thirteen components in the specification list and
nothing else}. Absent from the model, and therefore from the cost, are at least:
fuel cladding and assembly hardware, control and shutdown rods and their
drives, in-vessel and ex-vessel shielding, the secondary sodium circuit beyond
the IHX boundary, steam generators and the power conversion system, the
containment and reactor building, and all instrumentation, valves and piping.

This is appropriate at zero-level design --- the tool costs the primary-system
hardware it actually draws --- but the resulting figure is a
\emph{primary-system material bill}, not a plant capital cost, and it should not
be quoted as one. Anything not in \code{SPECS} is not in the report, and the
report has no way to know it is missing.

% =====================================================================
\section{Extending and adapting the reporter}
% =====================================================================

\begin{description}
  \item[Changing prices.] Edit \code{COSTS} in \file{cost\_report.py}, or pass a
    \code{costs=} dictionary to \code{write\_cost\_report()} without touching the
    module. The latter is preferable for scenario studies, since several price
    scenarios can then be run over one geometry in a single script.
  \item[Changing currency.] Set \code{currency=} to a matching label. Remember
    that this performs no conversion.
  \item[Alternative material libraries.] Pass \code{materials=} to price the
    same geometry against a different composition set.
  \item[Adding a material.] Add it to \file{materials.py} in either accepted
    form, and add a price to \code{COSTS}. No change to the reporter is needed.
  \item[Adding a component.] If it is single-material, nothing is required ---
    \code{generic\_zones} handles it. If its internals need more than one
    material, add a zone function to \code{MATERIAL\_ZONES}; the reporter picks
    it up automatically, as does the homogeniser. Mark any region that is
    contained rather than constructed with \code{role="fluid"} and it will be
    measured for neutronics but kept off the bill.
  \item[Programmatic access.] \code{material\_volumes(specs)} returns the raw
    \code{(obj\_id, obj\_type, zone, material, volume)} rows without writing any
    file, which is the natural hook for custom analyses, optimisation loops or
    alternative output formats.
\end{description}

% =====================================================================
\section{Reproducing the results in this document}
% =====================================================================

\begin{lstlisting}[basicstyle=\ttfamily\scriptsize]
arch -arm64 .venv/bin/python cost_report.py    # density + unit-cancellation self-test
arch -arm64 .venv/bin/python homogenise.py     # atom-conservation self-test
arch -arm64 .venv/bin/python examples_reactor/example_final_esfr_smr_with_materials.py
\end{lstlisting}

\code{arch -arm64} is required on this machine: the installed OCP/CadQuery
wheels are arm64, and a shell running under Rosetta cannot load them.

% =====================================================================
\section{Summary}
% =====================================================================

The cost reporter is a thin, deliberately un-clever layer over machinery that
already exists for neutronics. Its complete contribution is: a price list in
currency per kilogram, one unit-conversion constant per supported length unit,
a density derivation from number densities (Equation~\eqref{eq:rho}), a
solid/fluid scope rule, and a spreadsheet writer. Everything else --- what the
parts are, how big they are, what they are made of --- is read from the single
specification list that also produces the CAD assembly and the homogenised
neutronics model. The consequence is that the bill of materials cannot become
stale relative to the design, and that any change to the geometry is reflected
in the cost on the next run at the price of roughly one additional assembly
build.

The corollary, made concrete by Section~\ref{sec:audit}, is that the report is
exactly as good as the specification it reads. It cannot detect that a fuel
composition has been applied to four times the height it was homogenised over;
it can only price it accurately.

\end{document}
